\documentclass[11pt]{article}
\usepackage{campaign}

\title{Model Bill: The Capacity-Based Medical Aid in Dying Act (US)}
\author{Justin Bogner}
\date{2026}

\begin{document}

\maketitle

\begin{abstract}
\noindent \textit{Note to Legislative Counsel: This model bill is designed to replace or amend existing MAiD statutes (which rely on terminal prognoses) with a strict decisional-capacity standard. It incorporates the ``structural accountability mechanism'' reporting requirements.}
\end{abstract}

\section*{SECTION 1. SHORT TITLE.}
This Act may be cited as the ``Capacity-Based Medical Aid in Dying Act.''

\section*{SECTION 2. LEGISLATIVE INTENT.}
The Legislature finds and declares that:
\begin{enumerate}
    \item[(a)] The right to bodily autonomy includes the right of a competent adult to choose a peaceful death.
    \item[(b)] Current statutes restricting Medical Aid in Dying (MAiD) to individuals with a terminal prognosis of six months or less are arbitrary and philosophically inconsistent with the established right to refuse life-sustaining medical treatment.
    \item[(c)] Eligibility for MAiD must be based on a rigorous assessment of an individual’s decisional capacity, rather than a subjective assessment of suffering or life expectancy.
    \item[(d)] The state has a compelling interest in ensuring that requests for MAiD are not driven by systemic failures to provide adequate healthcare, housing, or disability support. Therefore, mandatory socio-economic data collection is required to hold the state accountable for the conditions that precipitate MAiD requests.
\end{enumerate}

\section*{SECTION 3. DEFINITIONS.}
As used in this Act:
\begin{itemize}
    \item \textbf{(a) ``Attending Physician''} means the physician who has primary responsibility for the care of the patient and treatment of the patient's condition.
    \item \textbf{(b) ``decisional capacity''} means the ability to understand and appreciate the nature and consequences of a health care decision, including the significant benefits and harms of and reasonable alternatives to any proposed health care.
    \item \textbf{(c) ``Independent Psychiatric Evaluator''} means a licensed psychiatrist who is not the attending physician and is not related to the patient by blood or marriage.
\end{itemize}

\section*{SECTION 4. ELIGIBILITY.}
An adult (18 years of age or older) who is a resident of this state shall be eligible to request and receive Medical Aid in Dying if they meet the following criteria:
\begin{enumerate}
    \item They have made a voluntary, written request for MAiD.
    \item They have undergone the two-stage capacity assessment as defined in Section 5 and have been found to possess durable decisional capacity.
\end{enumerate}

\section*{SECTION 5. TWO-STAGE CAPACITY ASSESSMENT.}
\begin{itemize}
    \item \textbf{(a) Stage One (Initial Assessment):} The attending physician must determine that the patient possesses decisional capacity and that the request is voluntary and not the result of acute coercion.
    \item \textbf{(b) Waiting Period:} Following the Stage One assessment, there shall be a mandatory waiting period of no less than 90 days for applicants whose natural death is not reasonably foreseeable. For applicants with a terminal prognosis, no mandatory waiting period is required, and provision may occur as soon as is feasibly possible.
    \item \textbf{(c) Stage Two (Independent Psychiatric Review):} Following the waiting period, the patient must be evaluated by an Independent Psychiatric Evaluator. 
    \begin{itemize}
        \item[(1)] The Evaluator must certify that the patient possesses durable decisional capacity.
        \item[(2)] A diagnosis of clinical depression or other mental health disorder shall \textit{not} automatically disqualify a patient. The Evaluator must determine if the disorder actively impairs the patient's ability to make a rational, informed, and durable decision regarding MAiD. If it does not impair capacity, the patient remains eligible.
    \end{itemize}
\end{itemize}

\section*{SECTION 6. THE STRUCTURAL ACCOUNTABILITY MECHANISM (SOCIOECONOMIC REPORTING).}
\begin{itemize}
    \item \textbf{(a)} Prior to the prescription of MAiD medication, the attending physician must conduct a standardized interview regarding the primary drivers of the patient's request.
    \item \textbf{(b)} The attending physician must report to the State Department of Health the extent to which the following factors influenced the patient's decision:
    \begin{itemize}
        \item[(1)] Inability to access or afford necessary medical care, palliative care, or psychiatric care.
        \item[(2)] Inability to access or afford adequate housing.
        \item[(3)] Lack of sufficient disability support services or home-care assistance.
    \end{itemize}
    \item \textbf{(c)} The State Department of Health shall compile this data into an annual, anonymized public report. 
    \item \textbf{(d)} Under no circumstances shall the presence of socioeconomic drivers invalidate a capacitated patient's eligibility for MAiD. The data is collected solely for the purpose of state accountability and policy reform.
\end{itemize}

\section*{SECTION 7. IMMUNITY.}
No physician, pharmacist, or healthcare provider who participates in good faith compliance with this Act shall be subject to civil liability, criminal prosecution, or professional disciplinary action.

\end{document}
